\documentclass[a5paper,10pt]{article}

% https://github.com/InDevRus/filippov-solutions

% Нумерация страниц
\pagenumbering{gobble}

% Настройка полей
\usepackage{geometry}
\geometry{left=1cm}
\geometry{right=1cm}
\geometry{top=1cm}
\geometry{bottom=1.5cm}

% Вставка таблицы
\usepackage{array}
\usepackage{tabularx}

% Инструменты выравнивания формул
\usepackage{amsmath}
\usepackage{mathtools}

% Диакритика в математике
\usepackage{amsfonts}

% Вычисления размеров на ходу
\usepackage{calc}

% Удобные записи производных
\usepackage[italic=true]{derivative}

% Настройки сносок
\usepackage{enotez}

% Длинные знаки векторов
\usepackage{anyfontsize}
\usepackage[b]{esvect}

% Вставка картинок
\usepackage{graphicx}

% Вставка красной строки
\usepackage{indentfirst}
\setlength{\parindent}{1em}

% Условия в командах
\usepackage{xifthen}

% Луа-скрипты
\usepackage{luacode}

% Настройка команд отдельно в математическом и текстовом режимах
\usepackage{mathcommand}

% Для повторения LaTeX кода
\usepackage{multido}

% Конфигурация отступов формул
\delimitershortfall-1sp
\usepackage{mleftright}
\mleftright

% Растягивание объектов
\usepackage{relsize}

% Обтекание текстом
\usepackage{wrapfig}
\usepackage{subcaption}
\usepackage{floatflt}

% Поддержка ввода кириллицы
\usepackage[english, russian]{babel}

% Настройка корневой позиции для компилятора
\usepackage{import}
\providecommand*{\ProjectRootPath}{..}

\import{\ProjectRootPath/preambles/macros/}{theorems}

\import{\ProjectRootPath/preambles/fonts}{font_setup}

% Знак градуса
\usepackage{siunitx}

\import{\ProjectRootPath/preambles/macros/}{operators}
\import{\ProjectRootPath/preambles/macros/}{redefinitions}

\import{\ProjectRootPath/preambles/fonts}{letters_setup}
\import{\ProjectRootPath/preambles/fonts/adjustments}{custom_superscripts}

\import{\ProjectRootPath/preambles/custom}{formatting}
\import{\ProjectRootPath/preambles/custom}{frames}
\import{\ProjectRootPath/preambles/custom}{list_setup}

% TODO: перевести команды на современный стиль объявления

\begin{document}

Уравнение $y\`\`\br{x} - 2\br{x - 1} y\`\br{x} + \pow{x}{2} y = 0$
необходимо привести к каноническому виду прежде, чем применять 
преобразование Лиувилля (задача \textbf{737}). Это можно сделать,
например, заменой
$$y\br{x} = \exp \br{- \frac {1} {2} \tint {p\br{x} \di x} }\, z\br{x}.$$
В нашем случае 
$p\br{x} = -2\br{x - 1}$, $q\br{x} = \pow{x}{2},$
так что заменим
$$y\br{x} = \exp\br{\frac {\br{x - 1}^2} {2}}\, z\br{x}.$$
Инвариант уравнения равен
$$\sub{Q}{z\br{x}\!} \br{x} 
    = q\br{x} - \frac {1} {2} p\`\br{x} - \frac {1} {4} \pow{p}{2}\br{x}
    = 2x.$$
После замены имеем уравнение $z\`\`\br{x} + 2x z\br{x} = 0$. 
Теперь преобразуем полученное уравнение по Лиувиллю. 
В обозначениях задачи \textbf{737} 
$\psi\br{x} = \tpow {\br{2x}} {\elevate{-\frac {1} {4}}}$,
$$\phi\br{x} 
    = \tintlim{0}{x} {\br{\psi\br{s}}^{-2}} \di s 
    = \frac {2\sqrt{2}} {3} \tpow{x}{\elevate{\frac {3} {2}}}$$
Соответственно, необходимо заменить переменные следующим образом:
$$
\tsys{
    z = \psi\br{x} u 
        = \tpow {\br{2x}} {-\frac {1} {4}} u \\
    t = \phi\br{x} 
        = \dfrac {2\sqrt{2}} {3} 
        \tpow {x} {\elevate{\frac {3} {2}}}.
}
$$
Инвариант преобразованного в терминах $z\br{t}$ уравнения принимает вид
$$\sub{Q}{z\br{t}\!}\br{t} 
    = 1 + \pow{\psi}{3}\br{x} \psi\`\`\br{x} 
    = 1 + \frac {5} {34} \pow{x}{-3} 
    = 1 + O\br{\pow{x}{-3}} 
    = 1 + O\br{\pow{t}{-2}}$$
Окончательно, для фундаментальной системы решений уравнения 
$z\`\`\br{t} + \sub{Q}{z\br{t}} \,z\br{t} = 0$ получаем следующие оценки:
\begin{align*}
    \sub{z}{1}\br{t} = \cos\br{t} + O\br{\pow{t}{-1}}
    && \sub{z}{2}\br{t} = \sin\br{t} + O\br{\pow{t}{-1}}
\end{align*}
Подставляя обратно, получаем \\
\begin{align*}
    \sub{y}{1}\br{x} 
    & = \exp\br{\frac {\br{x - 1}^2} {2}} 
        \tpow{\br{2x}}{\elevate{-\frac {1} {4}}} \sub{z}{1}\br{t} \\
    & = \exp\br{\frac {\br{x - 1}^2} {2}} 
        \tpow{\br{2x}}{\elevate{-\frac {1} {4}}} 
        \br{\cos\br{\frac {2\sqrt{2}} {3} 
            \tpow{x}{\elevate{\frac {3} {2}}}} 
            + O\br{\tpow {t} {\elevate{-\frac {3} {2}}}}} \\ 
    & = \exp\br{\frac {\br{x - 1}^2} {2}} 
        \br{ \tpow {\br{2x}}{\elevate{-\frac {1} {4}}}
            \cos\br{\frac {2\sqrt{2}} {3} 
            \tpow{x}{\elevate{\frac {3} {2}}}}
            + O\br{\tpow{x} {\elevate{-\frac {7} {4}}}}},
\end{align*}
и аналогично для $\sub{y}{2}$.

\pagebreak
Наконец, оценка фундаментальной системы решений уравнения 
$$y\`\`\br{x} - 2\br{x - 1} y\`\br{x} + \pow{x}{2} y = 0$$
такова:
\begin{align*}
    \sub{y}{1}\br{x} 
    = \exp\br{\frac {\br{x - 1}^2} {2}} \br{ 
        \tpow{\br{2x}} {\elevate{-\frac {1} {4}}} 
        \cos\br{\frac {2\sqrt{2}} {3} 
        \tpow{x} {\elevate{\frac {3} {2}}}}
        + O\br{ \tpow{x} {-\frac {7} {4}}}}; \\
    \sub{y}{2}\br{x}
    = \exp\br{\frac {\br{x - 1}^2} {2}} \br{ 
        \tpow{\br{2x}} {\elevate{-\frac {1} {4}}}
        \sin\br{\frac {2\sqrt{2}} {3} 
        \tpow{x} {\elevate{\frac {3} {2}}}}
        + O\br{ \tpow{x} {-\frac {7} {4}}}}.
\end{align*}

% TODO: добавить остальные

\end{document}
